\subsection{Stochastic stepping trace estimation}
\label{sec:appendix_sste_method}

SSTE has been described in \cref{subsec:randomized_trace_estimation} and is not repeated here. Let
\(N_p\) be the number of probes used for the estimation, and let the cost account for the adjoint
stepper as well, so that
\begin{equation}
  \mathcal{C}_{\mathrm{SSTE}}(T)=6N_p
  \label{eq:sste_fixed_time_cost}
\end{equation}
for a fixed time. The real advantage at a single target time is then guaranteed whenever
\(\mathcal{C}_{\mathrm{SSTE}}\ll\mathcal{C}_{\mathrm{DCT}}\), that is, whenever \(N_p\ll r/6\). We
expect this in problems whose dynamics are low rank, because the sketch captures those directions
through the action on random probes, and the examples reported below confirm it.

One major bottleneck nevertheless remains for an evaluation across \(N_t\) intervals, where the
cost grows as \(O(N_t^{2})\), because the backward loop must be performed several times over the
same instants in time, as \cref{fig:sste_time_history_cost} reports. The cost is then
\begin{equation}
  \sum_{j=1}^{N_t} j
  =
  \frac{N_t(N_t+1)}{2}
  \sim
  \frac{1}{2}N_t^{2},
  \label{eq:sste_backward_sum}
\end{equation}
which brings the total stepping work to
\begin{equation}
  \mathcal{W}_{\mathrm{SSTE}}=3N_pN_t^{2}.
  \label{eq:sste_multiple_time_cost}
\end{equation}

\begin{figure}[t]
  \centering
  \begin{tikzpicture}[x=1.15cm, y=1cm, >=Stealth, line width=0.7pt]
    % one row per target time: the forward sweep is stored once, while the
    % adjoint loop restarts from every target time and repeats the same intervals
    \foreach \row/\nback/\col in {0/1/blue, 1/2/green!55!black, 2/3/magenta}{%
      \begin{scope}[yshift={-\row*2.0cm}]
        \draw[black] (-0.3,0) -- (7.3,0);
        \foreach \i in {0,...,7}{\fill (\i,0) circle (1.5pt);}
        \foreach \i in {0,...,6}{\draw[orange,->] (\i,0) to[bend left=55] (\i+1,0);}
        \node[anchor=south, font=\footnotesize] at (3.5,0.62)
          {forward: solved once, \(\hat{\mathbf{q}}(t)\) stored};
        \node[anchor=west, font=\footnotesize] at (7.6,0) {\(T_{\the\numexpr\row+1\relax}\)};
      \end{scope}
    }
    \foreach \j in {1}{\draw[blue,->] (\j,-0.1) to[bend left=55] (\j-1,-0.1);}
    \begin{scope}[yshift=-2.0cm]
      \foreach \j in {1,2}{\draw[green!55!black,->] (\j,-0.1) to[bend left=55] (\j-1,-0.1);}
    \end{scope}
    \begin{scope}[yshift=-4.0cm]
      \foreach \j in {1,2,3}{\draw[magenta,->] (\j,-0.1) to[bend left=55] (\j-1,-0.1);}
    \end{scope}
  \end{tikzpicture}
  \caption{Stepping work of SSTE over a history of target times. The forward solve is performed
  once and stored, whereas the adjoint loop restarts from each target time \(T_j\) and traverses
  the same intervals again, so that the \(j\)-th target time costs \(j\) backward steps and the
  total work grows as \(O(N_t^{2})\).}
  \label{fig:sste_time_history_cost}
\end{figure}

The method must therefore always be preferred whenever
\(\mathcal{W}_{\mathrm{SSTE}}\ll\mathcal{W}_{\mathrm{DCT}}\), that is, whenever
\(3N_pN_t^{2}\ll rN_t\), which reduces to \(3N_pN_t\ll r\). This can be useful only when
\(r\sim N\) and \(N_p\sim O(1)\). Adjoint methods are known to carry this bottleneck and, as the
structure of the problem shows, they are very amenable to parallelization in time, as discussed by
\citet{skene-eggl-schmid-2021}; we do not discuss such strategies here.
